\documentclass[12pt]{article}
\usepackage{amsmath, amssymb, amsthm}
\usepackage{geometry}
\geometry{margin=1in}

\newtheorem{exercise}{Exercise}
\theoremstyle{remark}
\newtheorem*{solution}{Solution}
\newtheorem*{remark}{Remark}

\title{Random Walks --- Week 8 Piazza Exercises}
\author{CUNY, Spring 2026}
\date{}

\begin{document}
\maketitle

\section*{Week 8: Martingales}

Let $(\Omega, \mathcal{F}, P)$ be a probability space with filtration $(\mathcal{F}_n)_{n \ge 0}$.
Recall that a sequence $(M_n, \mathcal{F}_n)$ is a \emph{martingale} if:
(i) $M_n$ is $\mathcal{F}_n$-measurable, (ii) $E|M_n| < \infty$, and
(iii) $E[M_{n+1} \mid \mathcal{F}_n] = M_n$ a.s.  It is a \emph{supermartingale}
(resp.\ \emph{submartingale}) if (iii) is replaced by $\le$ (resp.\ $\ge$).

% ----------------------------------------------------------------
\begin{exercise}[Conditional expectation is a martingale]
Let $(M_n, \mathcal{F}_n)$ be a super- (sub-)martingale.
\begin{enumerate}
  \item[(a)] Show that $E[M_n \mid \mathcal{F}_n] = M_n$ a.s.\ for all $n$.
  \item[(b)] Show that $(M_n)$ is a martingale if and only if
             $E[M_{n+k} \mid \mathcal{F}_n] = M_n$ a.s.\ for all $n, k \ge 0$.
\end{enumerate}
\end{exercise}

\begin{solution}
\textbf{(a)} Since $M_n$ is $\mathcal{F}_n$-measurable, by the property
``taking out what is known,''
\[
  E[M_n \mid \mathcal{F}_n] = M_n \quad \text{a.s.}
\]

\textbf{(b)} $(\Rightarrow)$: If $(M_n)$ is a martingale, then by repeated conditioning
and the tower property,
\[
  E[M_{n+k} \mid \mathcal{F}_n]
  = E\bigl[E[M_{n+k} \mid \mathcal{F}_{n+k-1}] \mid \mathcal{F}_n\bigr]
  = E[M_{n+k-1} \mid \mathcal{F}_n]
  = \cdots = M_n.
\]
$(\Leftarrow)$: Taking $k = 1$ gives the martingale property directly.
\end{solution}

% ----------------------------------------------------------------
\begin{exercise}[Conditional expectation may take more values than the original]
Suppose $X$ takes exactly $n$ distinct values on $(\Omega, \mathcal{F}_0, P)$, and
$\mathcal{F} \subseteq \mathcal{F}_0$ is a sub-$\sigma$-algebra.  Is it necessarily true
that $E(X \mid \mathcal{F})$ takes at most $n$ distinct values?
\end{exercise}

\begin{solution}
\textbf{No.}  Here is a counterexample with $n = 2$.

Let $\Omega = [0,1)$ with Lebesgue measure $P$.  Define the ``finer''
$\sigma$-algebra $\mathcal{F}_0$ to be generated by the four intervals
\[
  I_1 = [0,\tfrac{1}{4}),\quad I_2 = [\tfrac{1}{4},\tfrac{1}{2}),\quad
  I_3 = [\tfrac{1}{2},\tfrac{3}{4}),\quad I_4 = [\tfrac{3}{4},1),
\]
and the ``coarser'' sub-$\sigma$-algebra $\mathcal{F}$ generated by
\[
  A_1 = [0,\tfrac{1}{4}),\quad A_2 = [\tfrac{1}{4},\tfrac{3}{4}),\quad A_3 = [\tfrac{3}{4},1).
\]

Define
\[
  X(\omega) = \begin{cases} 0 & \text{if } \omega < \tfrac{1}{2}, \\ 1 & \text{if } \omega \ge \tfrac{1}{2}. \end{cases}
\]
So $X$ takes exactly $n = 2$ distinct values, $\{0, 1\}$.

Now compute $E(X \mid \mathcal{F})$.  On each atom of $\mathcal{F}$:
\begin{align*}
  E(X \mid \mathcal{F})(\omega) &= \frac{E(X \cdot \mathbf{1}_{A_i})}{P(A_i)} \quad \text{for } \omega \in A_i.
\end{align*}
\begin{itemize}
  \item On $A_1 = [0,\tfrac{1}{4})$: $X = 0$ throughout, so $E(X \mid \mathcal{F}) = 0$.
  \item On $A_2 = [\tfrac{1}{4},\tfrac{3}{4})$: $X = 0$ on $[\tfrac{1}{4},\tfrac{1}{2})$ and $X = 1$ on $[\tfrac{1}{2},\tfrac{3}{4})$, each of length $\tfrac{1}{4}$, so
        $E(X \cdot \mathbf{1}_{A_2}) = \tfrac{1}{4}$ and $P(A_2) = \tfrac{1}{2}$, giving $E(X \mid \mathcal{F}) = \tfrac{1}{2}$.
  \item On $A_3 = [\tfrac{3}{4},1)$: $X = 1$ throughout, so $E(X \mid \mathcal{F}) = 1$.
\end{itemize}

Therefore
\[
  E(X \mid \mathcal{F})(\omega) = \begin{cases}
    0 & \omega \in [0,\tfrac{1}{4}), \\
    \tfrac{1}{2} & \omega \in [\tfrac{1}{4},\tfrac{3}{4}), \\
    1 & \omega \in [\tfrac{3}{4},1).
  \end{cases}
\]
This takes \textbf{3} distinct values $\{0, \tfrac{1}{2}, 1\}$, which is strictly more
than $n = 2$. \hfill$\square$

\begin{remark}
Intuitively, conditioning on a \emph{coarser} $\sigma$-algebra ``averages'' $X$ over
larger sets, which can create new intermediate values not originally in the range of $X$.
\end{remark}
\end{solution}

\end{document}
